\documentclass[10pt]{amsart}
\usepackage{graphicx,wrapfig,amsmath,amsthm,amssymb,latexsym}%,wrapfig,floatflt}
%\usepackage[matrix,arrow,cmtip,curve,dvips]{xy}
\DeclareGraphicsExtensions{.png}
\DeclareMathOperator{\im}{im}
\newcommand{\id}[1][{}]{\mathrm{Id}_{#1}}
\newcommand{\Q}{\mathbf Q}      % adds blackboard bold macros for
\newcommand{\R}{\mathbf R}      % conventional notation for
\newcommand{\Z}{\mathbf Z}      % rationals, reals, integers,
\newcommand{\C}{\mathbf C}      % complexes, naturals, field of p
\newcommand{\N}{\mathbf N}      % elements, p-adic integers
\newcommand{\D}{\mathbf D}
\newcommand{\Zp}{\mathbb{Z}_p}      % and numbers.
\newcommand{\Qp}{\mathbb{Q}_p}
\newcommand{\Fp}{\mathbb{F}_p}
\newcommand{\ie}{\emph{i.e.}}
\newcommand{\eg}{\emph{e.g.}}
\newcommand{\ind}[2]{[#1:#2]}       % index of groups, degree of field
\newcommand{\aut}[2][{}]{\mathrm{Aut}_{#1}(#2)}  % automorphism group
\newcommand{\isom}{\cong}       % isomorphic
\newcommand{\fk}[1]{\mathfrak{#1}}  % fraktur shorthand
\newcommand{\abs}[2][{}]{|#2|_{#1}} % generalized absolute value
\newcommand{\ord}[2][{}]{v_{#1}(#2)}    % roman ord(x), ord_p(x)
\newcommand{\Ex}{\begin{ex}}
\newcommand{\Eex}{\end{ex}}
\newcommand{\tr}[1]{#1^t}
\newcommand{\mat}[2]{M_{#1}(#2)}
\newcommand{\GL}[2]{\mathrm{GL}_{#1}(#2)}
\newcommand{\Lie}[1]{\mathrm{Lie}(#1)}
\newcommand{\ddt}[1]{\left. \frac{\mathrm{d}}{\mathrm{d}t} \right|_{t=#1}}
\newcommand{\ten}{\otimes}
\newcommand{\mg}[1]{{#1}^{\times}}
\newcommand{\wt}[1]{\widetilde{#1}}
\newcommand{\Tor}[2]{\mathrm{Tor}^{#1}_{#2}}

\theoremstyle{plain}            % define numbered and naked theorem
\newtheorem{ntheorem}{Theorem}[section]     % environments, numbered
\newtheorem{nlemma}[ntheorem]{Lemma}        % consecutively within
\newtheorem{nprop}[ntheorem]{Proposition}   % chapters
\newtheorem{ncor}[ntheorem]{Corollary}
\newtheorem*{theorem}{Theorem}
\newtheorem*{lemma}{Lemma}
\newtheorem*{prop}{Proposition}
\newtheorem*{cor}{Corollary}

% \newcounter{ntheorem}

\theoremstyle{definition}
\newtheorem{ndefn}[ntheorem]{Definition}
\newtheorem{nex}[ntheorem]{Example}
\newtheorem*{defn}{Definition}
\newtheorem*{ex}{Example}

%\theoremstyle{remark}
\newtheorem*{rmk}{Remark}
\newtheorem*{ntn}{Notation}

\usepackage[total={7.0in,10.0in},left=1.0in,right=1.0in]{geometry}
\setlength{\parskip}{7.5pt} \setlength{\parindent}{0pt}
%\pretolerance=4000 \setlength{\topmargin}{-1.0in}
%\setlength{\textheight}{10.0in} \setlength{\textwidth}{7in}
%\setlength{\headheight}{26pt} \setlength{\headsep}{8pt}
%\setlength{\oddsidemargin}{-0.25in}
%\setlength{\evensidemargin}{-0.25in}
\title{{\Large Mathematics 251, Fall 2011 \\ September 2}}
\pagestyle{empty}
\begin{document}
\maketitle
\thispagestyle{empty}

\begin{wrapfigure}{r}{140px}
    \begin{tabular}{ccccc} %%%%% set heights to 60px for two figures
	   \hspace*{-5mm} & \includegraphics[height=55px]{qr-facpage.png} & \hspace*{0mm} &
        \includegraphics[height=55px]{qr-homework.png} & \hfill
    \end{tabular}
%    \caption{web}
\end{wrapfigure}

\textbf{Instructor:} Dr.\ Dave Rosoff  \\
\textbf{Office:} Boone Hall 102C \\
\textbf{Office hours:} M 11--12, T 3:30--4:30, W 9:30--10:30, $\mathrm{\Theta}$ 1:30--2:30 \hspace*{0.25in}\\
\textbf{Email:} \verb[drosoff@collegeofidaho.edu[ 
\begin{center}

{\large \emph{The pursuit of knowledge, brother, is the askin' of many questions.}\footnote{Raymond Chandler, \emph{Farewell, My Lovely}.}}

\end{center}
\textbf{Text:} The text is \emph{Calculus: Early Transcendentals} by Jon Rogawski, \emph{first} edition. Although the second edition is available, many of you will already own copies of the first edition, and the department has elected to retain it for this course. It is a monolithic calculus book of a very standard form. We will assume the contents of Chapters 1--10; we'll cover chapters 11--16 and as much of 17 as possible, with some omissions.

\textbf{Course objectives:} Successful students in Math 251 will demonstrate mastery of univariate calculus sufficient to investigate vector analysis including Frenet--Serret (TNB) frames, parametric equations, elementary partial differentiation, and integration over rectangles and parallelepipeds using Fubini's theorem. They will use the change of variable theorem to apply these ideas to more general regions of integration (of full dimension) and finally, if time permits, develop the elementary theory of line and surface integrals with an eye toward Green's theorem, Stokes's theorem, and Gauss's divergence theorem.

\textbf{Course overview:} We generalize the main ideas and results of single-variable calculus to multivariate situations. There are two ways to proceed. First, we keep the domains of our functions the usual real numbers (or appropriate subsets thereof) and let their codomains (ranges) be higher-dimensional Euclidean spaces. This approach entails an investigation into the related concepts of \emph{parametrization} and \emph{vectors}. We may also ask what happens in the reverse scenario, when the domain of the functions at hand is allowed to be 2-, 3-, or higher-dimensional, and the function's values are real numbers in the usual sense. Here we will meet the essential concepts of \emph{partial derivative} and \emph{total} or \emph{Jacobian derivative} and revisit the familiar themes of differential calculus (related rates, optimization, etc.). Just as functions of several variables may be differentiated via their partial derivatives, so too can they be integrated over appropriate \emph{regions} (rather than intervals) in their domains. We study the important change-of-variable theorem that allows integrals over complicated regions to be reckoned in terms of simpler ones (e.g., rectangles) and examine some applications.

Many students encounter a significant increase in conceptual difficulty on passage to the third course in calculus. This is because of the multivariate nature of the investigations. There are notions in this realm that merge when specializing to the one-dimensional case: for example, a $1$-vector is the same thing as a point in $\mathbf{R}$, which is to say \emph{a number}. One must be careful when identifying $n$-vectors with points in $\mathbf{R}^n$, even though such identification is  legitimate. Another way to understand this phenomenon is that familiar ideas or objects may ramify, on passage to higher dimensions, into several parts: so that, for example, there are many types of $2$-dimensional analogues of the $1$-dimensional object we call ``open interval''.

The last part of the course comprises first steps toward a synthesis of the two approaches mentioned above. The technology of \emph{vector fields} is employed to study integrals not over flat rectangles, but over curves and (curved) surfaces. This is the very beginning of the important field of \emph{differential geometry}. The laws of electromagnetism, Maxwell's equations, are formulated in this language. We will make as much progress as we can toward the fundamental theorems of G.\ Green, C.\ F.\ Gauss, and G.\ Stokes\footnote{Famously, Stokes's theorem is due not to Stokes, but to Lord Kelvin; Stokes merely \emph{assigned} it, in 1854.}. All three are vast generalizations of the usual Fundamental Theorem of Calculus, and all have extremely important implications for physics, engineering, and the rest of mathematics.

\newpage

\textbf{Homework:} Homework in this class comprises both transcription and traditional written assignments.
\begin{itemize}
    \item Daily assignments from the text to be transcribed \emph{by hand} will be assigned and due at the beginning (8 am) of the next lecture. These exercises are to encourage reading of the text in advance of the lecture. Student learning improves when you arrive to class with some familiarity with the major terminology and concepts.
    \item Daily written assignments, due the following lecture day. See below for writing guidelines.
    \item Occasional lab assignments using Mathematica or group exercises assigned as homework.
\end{itemize}
Whenever you are writing a solution to a math problem, it is important to strive for the clearest exposition you can manage. Good mathematical writing is essential for anyone who wishes to think clearly about mathematics---sloppy writing invariably reflects underlying sloppy thinking. The process of making your ideas and reasoning \emph{clear, complete, and unambiguously correct} is the most powerful amplifier of mathematical power there is. Hence your solutions should be composed in brilliant English prose (e.g., accepted scientific usage, more or less correct grammar and spelling, and above all \emph{complete sentences}) sprinkled with tangy, delicious equations here and there. Solutions in the popular ``pile-of-equations'' style are to be avoided and will not get much credit. You must explain what is happening as the action unfolds.

I encourage all of you to form study groups and collaborate on your homework; each student is of course individually responsible for their own work. Collaborators must be acknowledged. \textbf{Late homework is generally not accepted without significant penalties}, and its acceptance is determined on a case-by-case basis according to the student's situation. No written homework will be accepted after it is returned to the class.

\textbf{Quizzes:} Unannounced quizzes will be given intermittently throughout the term. Some may be open-book and some may be group exercises. Missed quizzes are scored zero unless the absence from class is arranged well in advance, and occurs for a legitimate and compelling reason. In case of illness, make-ups may be administered; \textbf{contact Dr.\ Rosoff immediately} if you are sick and unable to come to class.

\textbf{Exams:} Three exams are given in class (see below for dates). A missed exam results in an exam grade of zero. Make-up policies are similar to those outlined above for quizzes. Arrangements for absences, again, must be made well in advance (two weeks suffices).
\begin{itemize}
	\item Exam 1 (tentative): Wednesday, September 22
	\item Exam 2 (tentative): Wednesday, October 26
    \item Exam 3 (tentative): Tuesday, November 22
	\item Final Exam: Monday, December 5, 1:30--4:30, Boone 134
\end{itemize}

\textbf{Grading:} Scores are computed as a weighted average, with the following weights: reading homework and group assignments $0.05 = 5\%$, quizzes $0.05 = 5\%$, assignments $0.15 = 15\%$, three in-class exams $0.55 = 55\%$, and final exam $0.20 = 20\%$. Observe that the weights sum to $1 = 100\%$. The exact determination of letter grades from these scores depends on the final distribution of scores in the class, but you can expect a C for earning 75\% of the points, a C+ for 80\%, a B-- for 83\%, and so on.

\textbf{Academic honesty:} Students are expected to complete all graded work in accordance with the College Honor Code. Cheating will not be tolerated. All parties involved in cheating will receive a final grade of F and such cases will be referred to the Vice President for Academic Affairs, who is a lot scarier than I am.

\textbf{A note on studying math:} By now you have studied enough mathematics to have learned something about how it is that the material passes through your shapely skull and into your soft, spongy brain. Nevertheless, you may find that this course is rather more difficult than your previous calculus courses. To really understand it, we will have to dig into subtle distinctions and nuances that no one has asked you to think about before. The reasons for this are outlined above: in the one-dimensional world, they simply do not signify. It is also more difficult to picture what is going on mentally, again owing to the presence of extra dimensions. Part of what I'm here to tell you is that while the material may seem wholly new and unfamiliar, the underlying principles of calculus (what do derivatives \emph{do}? what are integrals \emph{for}?) are immutable.

\textbf{Special accommodations:} Students who have documented disabilities as addressed by the Americans With Disabilities Act and who need any test or course materials to be furnished in an alternative format should notify me immediately (during the first week of class).  Reasonable efforts will be made to accommodate the needs of such students.

\begin{center}
\emph{ {\LARGE Good luck this semester!} }
\end{center}

\end{document}
